% \documentclass[12pt,paper=a4,ngerman]{scrreprt}
\documentclass[12pt,paper=a4,ngerman]{report}
\usepackage[a4paper,width=150mm,top=25mm,bottom=25mm]{geometry}

\usepackage{babel}
\usepackage[T1]{fontenc}
\usepackage[default]{raleway}
\usepackage{hyperref}
\usepackage{eurosym}
\usepackage{url}
\usepackage{enumitem} % Für benutzerdefinierte Listen
\usepackage{musicography}
\usepackage{relsize}

\usepackage[dvipsnames]{xcolor}
\definecolor{p_petrol}{HTML}{9EC2C0}
\definecolor{p_mint}{HTML}{D3F6E8}

\usepackage{fancyhdr}
\pagestyle{fancy}
\fancyhf{} % reset header and footer
\renewcommand{\headrulewidth}{0pt}
\fancyhead[C]{\color{p_petrol}Beitragsordnung - Plan b Chor Bonn e.V.}
\fancyfoot[C]{\thepage}

\renewcommand{\baselinestretch}{1.15}

\setcounter{secnumdepth}{0} \setcounter{tocdepth}{3}

\hypersetup{
pdfauthor = {plan-b-chor.de},
pdftitle = {Beitragsordnung Plan b Chor Bonn e.V.},
colorlinks,
linkcolor=p_petrol,
urlcolor=p_petrol,
}

\title{\textbf{Beitragsordnung}}
\author{Plan b Chor Bonn e.V.}
\date{19. Januar 2026}

\newcounter{subparagraphCounter} % Zähler für Unterpunkte

% Definition eines neuen Befehls für Unterparagraphen
\newcommand{\mysubparagraph}{%
    \noindent(\alph{subparagraphCounter}) % Gibt den Buchstaben (a, b, c, ...) aus
    \refstepcounter{subparagraphCounter} % Erhöht den Zähler für Unterparagraphen um 1
    \hspace{1em}
}

%%%

\newcounter{paragraphCounter} % Erstellt einen neuen Zähler namens paragraphCounter
\setcounter{paragraphCounter}{1}

% Definition eines neuen Befehls für Paragraphen
\newcommand{\myparagraph}{
    \noindent\S\ \theparagraphCounter % Gibt das Paragraphensymbol und den Zähler aus
    \label{par-\theparagraphCounter}
    \refstepcounter{paragraphCounter} % Erhöht den Zähler um 1
    \setcounter{subparagraphCounter}{0} % Setzt den Unterabschnittszähler zurück
    \hspace{1em}
}

% Anpassung des Aufzählungsstils für itemize
\setlist[itemize,1]{label=\myparagraph, before=\addtocounter{paragraphCounter}{-1}, left=0pt} % 1. Ebene
\setlist[itemize,2]{label=\mysubparagraph, left=7pt} % 2. Ebene

\begin{document}
\maketitle

\section{}

Diese Ordnung wurde am 19.01.2026 beschlossen und gilt ab dem Stichtag 01.02.2026. Sie
setzt alle vorher beschlossenen Beitragsregelungen außer Kraft.

\section{Allgemeines}

% Verwende eine itemize-Umgebung mit dem angepassten Zähler als Aufzählungszeichen
\begin{itemize}
    \item Diese Beitragsordnung ist nicht Bestandteil der Satzung. Beschlüsse über die Änderungen der Beitragsordnung gelten ab dem auf die Beschlussfassung folgenden Monat.
    \item Beim Austritt aus dem Verein erfolgt keine Rückerstattung bereits geleisteter Beiträge.
    \item Die Mitgliedsbeiträge sind die Haupteinnahmequelle des Vereins. Die Mitgliederversammlung beauftragt den Vorstand, diese Mittel im Sinne des Zweckes des Plan b Chor Bonn e.V. vollständig zu verwalten.
\end{itemize}

\section{Zahlungsweise und Fälligkeit}

\begin{itemize}
    \item Die festgesetzten Beträge sind zum ersten Werktag jedes Monats auf das Beitragskonto des Vereins zu entrichten. Durch Beschluss der Mitgliederversammlung kann auch ein anderer Termin festgelegt werden.
    \item Der Verein bittet darum, die entsprechenden Beiträge nach Möglichkeit auf freiwilliger Basis zu Beginn eines (Halb-) Jahres jeweils für das kommende (Halb-) Jahr vorab zu bezahlen. Bei Beendung der Mitgliedschaft wird der zu viel bezahlte Beitrag zurückgezahlt.
    \item Folgende Regelungen gelten im Falle von ausbleibenden Mitgliedsbeiträgen:
    
    \begin{itemize}
        \item Bei ausbleibender Zahlung erfolgt eine Ermahnung per E-Mail. Erfolgt nach zehn Werktagen keine Zahlung, dann folgt eine weitere Ermahnung per E-Mail mit erneuter Frist von fünf Werktagen.
        \item Darauf folgen im Abstand von fünf Werktagen Mahnungen. Für Mahnungen werden Mahngebühren i.H.v 10 Euro je ausgesprochener Mahnung erhoben. Bleibt innerhalb der nächsten sechs Monate die Beitragszahlung erneut aus, so erfolgen direkt Mahnungen.
        \item Bei ausbleibenden Zahlungen nach mindestens sechs Monaten ohne ausbleibende Zahlung startet der Vorgang erneut mit zwei Ermahnungen.
    \end{itemize}
\end{itemize}

\section{Beiträge}

\begin{itemize}
    \item Ehrenmitglieder sind von der Beitragszahlung befreit.
    \item Die Höhe des Beitrags richtet sich nach der Form der Mitgliedschaft.
    \begin{itemize}
        \item Der Beitrag für aktive Mitglieder liegt bei \EUR{30}.
        \item Der Beitrag für passive Mitglieder liegt bei \EUR{8}.
        \item Der ermäßigte Mindestbeitrag liegt bei \EUR{20}.
    \end{itemize}

    \item Ermäßigte Beitragsformen müssen beantragt werden. Der Anspruch auf die Ermäßigung ist mit entsprechenden Unterlagen nachzuweisen. Der geschäftsführende Vorstand entscheidet über die Einstufung im Rahmen der von der Beitragsordnung vorgegebenen Beiträge.

    \item Sofern ein ermäßigter Beitrag in Anspruch genommen wird, sind Änderungen der persönlichen finanziellen Verhältnisse dem Vorstand schnellstmöglich mitzuteilen.

    \item Die Beiträge stellen Mindestbeiträge dar. Der Beitrag kann auf freiwilliger Basis aufgestockt werden.

    \item Bei Aufnahme in den Plan b Chor Bonn e.V. wird eine Aufnahmegebühr von \EUR{20} fällig.
\end{itemize}

\section{Vereinskonto}

\begin{itemize}
    \item Der Mitgliedsbeitrag ist auf das folgende Vereinskonto zu entrichten:\vspace{1em}\\
    \textbf{Bank}: VR-Bank Altenburger Land / Deutsche Skatbank\\
    \textbf{IBAN}: DE47 8306 5408 0005 4440 04
\end{itemize}

\section{Beitragsbefreiung im Härtefall}

\begin{itemize}
    \item Bei materieller Bedürftigkeit können Mitglieder von der Zahlung des Beitrages nach \hyperref[par-7]{§§7 - 11} der Beitragsordnung befreit werden.
    \item Alle Antragstellenden sind verpflichtet, ihre Einkommensverhältnisse wahrheitsgemäß darzulegen. Dem Antrag sind Unterlagen in Kopie beizufügen, aus denen die wirtschaftlichen Verhältnisse nach \hyperref[par-14]{§14} hervorgehen.
    \item Eine Befreiung von Pflichten ist stets befristet. Höchstens jedoch auf die Dauer eines Geschäftsjahrs.
    \item Über die Befreiung entscheidet der Vorstand. Der Antrag auf Befreiung erfolgt schriftlich. Bei Ablehnung kann die betroffene Person beim Härtefallausschuss Berufung einlegen. Der Härtefallausschuss entscheidet endgültig.
    \item Der Vorstand kann die Annahme des Antrages von der Vorlage notwendiger Nachweise abhängig machen. Insbesondere kann er Anträge, die nicht innerhalb eines Monats nach Eingang des Antrages vervollständigt oder durch geeignete Nachweise belegt wurden, als unvollständig ablehnen. Es liegt in der Verantwortung der antragstellenden Person, notwendige Unterlagen und Nachweise nachzureichen.
    \item Auf die Befreiung besteht kein Rechtsanspruch. Sie wird im Rahmen und auf Grundlage der zur Verfügung stehenden Haushaltsmittel gewährt.
    \item Die folgenden Personengruppen können bei der Antragsbearbeitung berücksichtigt werden:
    \begin{itemize}
        \item Studierende und Auszubildende
        \item Berufstätige im Falle von ausbleibenden Lohnzahlungen
        \item Schüler:innen, die ihren Beitrag selbst zahlen
        \item Arbeitslose
        \item Empfänger:innen von Sozialhilfe
        \item Personen mit anderen nachweisbaren Einschränkungen
    \end{itemize}
    \item Über Anträge nach \hyperref[par-14]{§14} entscheidet der Vorstand nach billigem und pflichtgemäßem Ermessen.
    \item Die Bewilligung eines Antrages erfolgt schriftlich und umfasst mindestens die Dauer und den Grund der Befreiung.
    \item Als anzurechnendes Einkommen im Sinne dieser Ordnung zählen staatliche Leistungen, wie z.B. Sozialhilfe oder Kindergeld. Nicht angerechnet wird das Elterngeld. Voll angerechnet werden Unterhaltsleistungen von Dritten an den/die Antragstellenden, wie durch Eltern, geschieden oder getrenntlebende Ehegatten, Väter oder Mütter der mit im Haushalt der*des Antragstellenden lebenden Kinder oder durch andere Personen. Eigenes Einkommen von Haushaltsmitgliedern der antragstellenden Person wird angerechnet. Vom Bruttoentgelt aus unselbstständigen Arbeitsverhältnissen werden die tatsächlich geleisteten Steuern und Sozialversicherungsbeiträge abgezogen.
    \item Die/der Antragstellende hat in angemessenem Umfang zur Entlastung ihrer/seiner finanziellen Situation beizutragen. Hinderungsgründe für die Aufnahme einer Erwerbstätigkeit, insbesondere der Erziehung von Kindern unter 12 Jahren, werden vom Vorstand nach billigem und pflichtgemäßem Ermessen anerkannt.
    \item Die Antragsfrist für ein Jahr endet am letzten Werktag der Woche nach Beginn des nächsten Geschäftsjahres. Verspätet eingegangene Anträge können ohne inhaltliche Prüfung zurückgewiesen werden. Es gilt jeweils das Datum des Posteingangs. In Sonderfällen kann der Vorstand Abweichungen tolerieren.
    \item Für Neueinsteigende endet die Antragsfrist 4 Wochen nach Beitritt in den Chor, innerhalb dieses Jahres.
    \item Anträge von Berechtigten nach \hyperref[par-20]{§20}, die nicht innerhalb der Frist nach \hyperref[par-25]{§25} eingehen, können ausnahmsweise berücksichtigt werden.
    \item Gegen jede Form von Entscheidungen des Vorstands im Härtefall kann das Widerspruchsverfahren stattfinden. Widerspruchsbehörde ist der Härtefallausschuss. Zur Wahrung der Sicherheit personenbezogener Daten tagt der Härtefallausschuss unter Ausschluss der Öffentlichkeit.
    \item Der Widerspruch ist innerhalb eines Monats nach Bekanntgabe des Bescheids schriftlich oder zur Niederschrift beim Vorstand zu erklären.
    \item Für das Widerspruchsverfahren finden die \href{https://www.gesetze-im-internet.de/vwgo/BJNR000170960.html\#BJNR000170960BJNG001001308}{§§68 ff. der Verwaltungsgerichtsordnung} Anwendung.
    \item Gegen die Entscheidungen des Härtefallausschusses ist der Verwaltungsrechtsweg gegeben.
    \item Falls die Zahl der bewilligten Beitragsbefreiungsanträge nicht zur vollständigen Ausschöpfung der im Haushalt des Vereins dafür vorgesehenen Mittel führt, wird eine zweckgebundene Rücklage für die folgenden Haushaltsjahre gebildet.
\end{itemize}

\section{Schlussbestimmungen}

\begin{itemize}
    \item Änderungen dieser Ordnung bedürfen der einfachen Mehrheit der anwesenden stimmberechtigten Mitglieder der Mitgliederversammlung.
\end{itemize}

\end{document}
